 \textbf{L'enjeu :} Face à un problème informatique, il existe souvent plusieurs solutions. Comment savoir laquelle est la plus efficace ? On pourrait simplement chronométrer le temps d'exécution, mais il y a un problème majeur...

\section*{Pourquoi ne pas simplement utiliser un chronomètre ? }

Le temps d'exécution réel dépend énormément de la machine utilisée. Un même programme sera bien plus rapide sur un ordinateur de compétition que sur un vieux PC de bureau.

\begin{center}
\begin{tabular}{llc}
     \hline
     \textbf{Machine} & \textbf{Fréquence} & \textbf{Capacité de calcul (estimée)} \\ \hline
     Vieux PC (année 2000) & 0,5 GHz & $500\,000\,000$ op/sec \\
     PC de la salle NSI & 3,5 GHz & $3\,500\,000\,000$ op/sec \\
     Serveur de calcul (nombreux c\oe urs) & 3 GHz par c\oe ur & Des centaines de milliards \\ \hline
\end{tabular}
\end{center}

Pour comparer les algorithmes de manière juste, on ne compte pas les secondes, mais le nombre d'\motcle{opérations élémentaires}. C'est la \textbf{complexité théorique}.

\vspace{0.3cm}
\hrule
\vspace{0.3cm}

\section*{Les Opérations Élémentaires : les briques de base}

\begin{Definition}[Opération élémentaire]
C'est une instruction de base dont le temps d'exécution est fixe : il ne change pas, peu importe la taille des données.
\end{Definition}

 \textbf{Le catalogue des opérations :}
\begin{itemize}
    \item \textbf{L'affectation :} \texttt{x = 5}.
    \item \textbf{Les calculs :} \texttt{a + b}, \texttt{x * 2}, \texttt{i += 1}.
    \item \textbf{Les comparaisons :} \texttt{if score > 100:} ou \texttt{while i < 10:}.
    \item \textbf{L'accès mémoire :} Lire la valeur de \texttt{variable} ou \texttt{liste[0]}.
\end{itemize}

\newpage

\subsection*{S'entraîner à compter : 5 cas d'école }

\begin{lstlisting}[style=pythonstyle, language=python]
# 1. Le cas de base (Calcul direct)
a = 10              # 1 op (aff)
res = (a + 3) * 5   # 1 add + 1 mult + 1 aff = 3 op
# Total : 4 operations elementaires.

# 2. Parcours d'une liste fixe
ma_liste = [10, 20, 30]     # 1 op (aff)
for nombre in ma_liste:     # 3 iterations :
    print(nombre + 1)       # 1 add + 1 affichage = 2 op
# Total : 1 + (3 * 2) = 7 operations elementaires.

# 3. La boucle while (Prudence !)
note = 15                   # 1 op (aff)
while note < 18:            # 3 tests vrais + 1 faux = 4 op
    note += 1               # 3 * (1 add + 1 aff) = 6 op
# Total : 1 + 4 + 6 = 11 operations elementaires.

# 4. Double boucle simple (imbriquee)
points = 0                  # 1 op (aff)
for i in range(2):          # Repete 2 fois :
    for j in range(3):      #   Repete 3 fois :
        points += 1         #     (1 add + 1 aff) = 2 op
# Total : 1 + (2 * 3 * 2) = 13 operations elementaires.

# 5. Appel de fonction dans une boucle
def ajoute_neuf(x):
    for i in range(3):      # 3 tours
        x += 3              # 1 add + 1 aff = 2 op
    return x                # 1 op

total = 0                   # 1 op
for j in range(3):          # 3 fois :
    # Fonction (7 op) + 1 appel + 1 aff = 9 op
    total = ajoute_neuf(total) 
# Total : 1 + (3 * 9) = 28 operations elementaires.
\end{lstlisting}